\section{Prompt Contracts and Convention Sheet}
\label{app:prompts}
\label{sec:prompt_contracts}
Prompts are versioned implementation artifacts; the hashes fingerprint the prompt constants in the released code. Every v3 prompt places the fixed blocks (schema, rules, convention sheet, examples) first and the item-specific blocks last, so hosted endpoints can reuse cached prefixes; the content is not shortened.
\begin{enumerate}\setlength\itemsep{1pt}
\item \textbf{Template proposal and induction} (\texttt{007be4}, \texttt{a827ce}, rules \texttt{5b8f6e}). A question with typed slots plus parameterized Cypher that uses every slot as \cy{\$slot}; conventions C4 and C10; a requested difficulty tier; induction also lists the templates already covered and asks for structurally different ones.
\item \textbf{Question realization} (\texttt{513b36}). Rewrite in one sampled register (terse analyst, business manager, imperative, or context-first) without changing meaning; of several rewrites, the one least similar to already accepted questions is kept; every bound value verbatim; the writer rules below.
\item \textbf{Reference generation} (\texttt{f4af10}). Schema-visible constructs and observed directions only; \cy{RETURN DISTINCT} for set answers; exact equality for quoted values; required context columns; no writes; the convention sheet. The template's own Cypher, with the bound values substituted, is kept as a fallback when the generated reference is rejected, and as an agreement check otherwise.
\item \textbf{Checklist judge} (\texttt{68dd5f}, frozen before the build). Inputs: the full schema with symmetric types marked, the convention sheet, the question, the Cypher, a clause-by-clause reading the judge writes first, five result rows, and deterministic notes (unlisted patterns, variables dropped at a \cy{WITH}). The judge answers nine criteria (values, relationships, aggregation, grouping, comparisons, ranking, columns, no extra constraints, unambiguous) with yes, no, or n/a and evidence; code passes the item only if no criterion is ``no'' and ``unambiguous'' is ``yes''. Ties, truncation, emptiness, read-only safety, and syntax are left to code.
\item \textbf{Downstream Text2Cypher} (standard and governed prompts). Read-only Cypher only; schema-visible constructs and direction preservation; the governed prompt adds the generator's conventions.
\end{enumerate}

\subsection{Convention Sheet (v3.2)}
The same text is given to the question writer, the reference generator, both judges, the outside auditor, and the human annotators.
\par\smallskip
\begingroup
\fontsize{6.3}{7.3}\selectfont
\begin{verbatim}
Benchmark conventions (questions are written under these
    rules and read with them):
C1 Values. Text in quotation marks is an exact, case-
    sensitive value. Identifiers are written without
    quotation marks (for example "account
    4803651952544056659") and match the label's
    identifying property. A quoted number or boolean that
    names a value ("'73'", "'False'") means that value in
    the property's stored type. Every value the question
    states appears in the query, and the query filters on
    no literal value that the question does not state.
C2 Direction. A verb with a direction ("transferred to",
    "invested in", "owned by", "replied to", "applied
    for") fixes the relationship direction. A neutral verb
    ("connected to", "linked with", "related to",
    "associated with") means either direction; a question
    that intends one direction must use a directed verb.
    Symmetric relationship types (listed with the schema)
    are matched without a direction.
C3 Roles. A role noun ("investor", "owner", "member",
    "author") covers every label that can play that role
    in the schema unless the question names the label.
C4 Counting. "How many <entities>" counts distinct
    entities. "How many <transfers / links / relationships
    / events / times>" counts relationship or event
    instances. Sum, average, minimum, and maximum are
    taken over all matching instances, never over DISTINCT
    values, unless the question says "distinct" or
    "different".
C5 Sets. Answers are compared as sets of rows: duplicate
    rows and row order do not matter unless the question
    asks for an order ("in chronological order", "sorted
    by", "ranked"). DISTINCT in a query is never an error
    by itself.
C6 Columns. The query returns every attribute the question
    asks for, by name. Extra columns that identify or
    describe the returned entities (ids, names, types,
    status flags) are allowed. Extra filters, exclusions,
    or row limits that the question does not state are not
    allowed.
C7 Comparisons. "After", "before", "more than", "less
    than", "over", "under", "exceeds" are strict (>, <).
    "On or after", "since", "at least", "at most", "no
    more than", "up to", "until" are inclusive (>=, <=).
    "Between A and B" includes both ends. "In <year>" or
    "in <month> <year>" covers the whole period. A two-way
    split ("before and after", "above and below") may put
    the boundary value in either group.
C8 Paths. "Within N hops" means 1 to N hops; "exactly N
    hops" means N; "directly" means 1 hop. A reachability
    question states its hop bound, and the query uses
    exactly that bound.
C9 Temporal chains. "Then", "followed by", "later", or
    "subsequently" between steps require each later step's
    timestamp to be greater than the earlier step's.
C10 Self-matches. "Other", "another", "else", or
    "different" (in the sense of "not the same one")
    require the query to exclude the anchor entity, for
    example a <> b for self-transfers. Without such a
    word, the query may either allow or exclude self-
    matches; neither is an error.
C11 Ranking. "Top k by X" or "the k highest X" sorts by X
    descending and keeps k rows; "lowest", "bottom",
    "fewest" sort ascending. "Which <entity> has the most
    X" returns the maximum. The ranking is over the
    population the question names; ranking a single given
    entity against itself is degenerate and not allowed.
C12 Grouping. "For each X", "per X", or "by X" groups by X
    and returns X's identifier with each aggregate.
    Without such words an aggregate is one value over the
    whole population. Groups with zero matches are
    included only when the question says so ("including
    those with none").
C13 One reading. A question has exactly one reasonable
    reading under these conventions. It contains no
    unfilled placeholders such as {k} and no malformed
    values (such as a quoted "account with id 123"), it
    names any time window, threshold, k, or hop bound it
    relies on, and it asks only for facts the graph
    records (no "current" or historical status the schema
    does not store).
\end{verbatim}
\endgroup

\subsection{Checklist Judge Prompt}
\label{app:judgeprompt}
Both panel judges and every chat-model auditor receive the prompt below, frozen before the build. Angle brackets mark the slots filled for each item: the conventions slot holds the sheet above, and the schema slot the full profiled schema.
\par\smallskip
\begingroup
\fontsize{6.3}{7.3}\selectfont
\begin{verbatim}
[system]
You are a meticulous reviewer of
    natural-language-to-Cypher benchmark items. Return
    strict JSON only. No markdown or extra text.

[user]
Task: decide whether a benchmark item is correct. An item
    is a natural-language question and a reference Cypher
    query over the graph schema below. The item is correct
    when the query returns exactly what the question asks
    for, read under the conventions.

<conventions>

Out of scope (code checks these before you see the item;
    never answer "no" because of them): read-only safety,
    Cypher syntax, empty results, ties at an ORDER BY ...
    LIMIT cutoff, and truncation by a LIMIT without ORDER
    BY. A large LIMIT that only caps a list (the result
    has fewer rows than the LIMIT) is harmless.
    Categorical value lists constrain literals written in
    the query, not values that appear in result rows. Code
    checks that labels and properties exist and that
    simple MATCH patterns follow the schema, but it does
    not check patterns inside EXISTS / NOT EXISTS
    subqueries, untyped nodes, or whether a pattern can
    relate the entities the question means; you must check
    those.

How to review: first read the query clause by clause and
    write down what it returns: which variables each MATCH
    binds, what each WITH carries forward (a variable that
    is not carried loses its binding and matches anything
    if it is used again), which entities each mandatory
    MATCH drops, whether separate patterns are combined
    into a cross product of rows, and what one result row
    represents. Then compare that reading with the
    question.

Checklist. For each criterion, first write short evidence
    that quotes the question and the query, then answer
    "yes", "no", or "na" (na only when the criterion does
    not apply).
1. values: every entity, identifier, number, date, and
    quoted value in the question appears in the query with
    the same meaning, and the query filters on no literal
    value the question does not state (C1). A quoted
    phrase that is not itself a stored value (for example
    'account with id 123' or 'company id company') is a
    no.
2. relationships: relationship types, node labels, and
    directions express the question's meaning, and the
    pattern connects the entities the question relates
    (C2, C3). Every pattern, including patterns inside
    EXISTS / NOT EXISTS and after WITH, uses a
    relationship the schema lists between those labels in
    the listed direction (symmetric types excepted);
    variables used after a WITH were carried through it;
    independent sub-answers are not combined into a cross
    product of rows.
3. aggregation: the aggregate function and the counted
    unit match the question (C4). na if the question asks
    for no aggregate.
4. grouping: the grouping key and the per-group output
    match the question (C12). A question that asks for one
    overall value gets one value, not one per entity; a
    per-entity question is not collapsed into one value.
    na if there is neither grouping nor an overall
    aggregate.
5. comparisons: comparison operators, boundaries, time
    windows, and hop ranges match the question (C7, C8,
    C9). na if there are none.
6. ranking: the sort key, sort direction, and k match the
    question (C11). A ranking whose population is
    restricted to one given entity (for example "among
    accounts with id 123, which have the most ...") or
    whose sort key is the same for every row is a no. na
    if the question asks for no ranking or ordering.
7. columns: the query returns every attribute the question
    asks for (C6).
8. no_extra_constraints: the query adds no filter,
    exclusion, hop bound, or row limit that the question
    does not state and the conventions do not require (C6,
    C8, C10), and no mandatory MATCH drops entities the
    question asks for (for example, keeping only a
    person's comments that reply to comments when the
    question asks for all of their comments).
9. unambiguous: the question has one reasonable reading
    under the conventions, contains no unfilled
    placeholder, asks only for facts the graph stores (no
    "current" job, change history, or population the
    schema lacks), and its answer is meaningful to an
    analyst (C11, C13).

Output: one JSON object with exactly these keys:
{"query_reading": "walk through the query clause by clause
    (each MATCH pattern as start label, type, end label;
    each WITH and what it carries; each WHERE; the
    RETURN), then say what one result row represents and
    which rows are kept or dropped", "values":
    {"evidence": "...", "answer": "yes|no|na"},
    "relationships": {"evidence": "...", "answer": "..."},
    "aggregation": {...}, "grouping": {...},
    "comparisons": {...}, "ranking": {...}, "columns":
    {...}, "no_extra_constraints": {...}, "unambiguous":
    {...}, "clarified_question": "if unambiguous is no,
    one rewrite of the question that has a single reading
    and matches the query; otherwise an empty string",
    "difficulty": "easy|medium|hard"}

Graph schema:
<schema>

Item to review
Question:
<question>

Reference Cypher:
<cypher>

Result: <row_summary>
<rows>

Query features (from the parser): <features>
Validator notes: <notes>

Return the JSON object now.
\end{verbatim}
\endgroup
